\subsection{2.7 评估与同期模型对比}\label{evaluation-and-comparison-with-contemporaneous-models}

在本节中，我们将 MAI-Base-1 与现有的同期基础模型进行对比，使用我们在开发中（如第 2.3 节所讨论）所用的同一组 NLL 评估的一个子集。其中一个是基于我们内部训练栈（第 2.8 节）的源代码；其余则由厂商委托提供，因此我们有信心它们不在任何模型的训练集中。我们没有使用表 3 中描述的多语言基准，这些基准来自我们已从自身训练中排除的公开来源，因此无法与这些模型进行可靠的比较。

在本对比中，我们聚焦于预训练基础模型，并报告跨分词器不变的每字节比特数（bits-per-byte，BPB）数值。大多数模型 API 无法进行 NLL 评估，因此我们聚焦于托管在 huggingface.co 上的开放权重模型：DeepSeek v3.2（DeepSeek-AI et al., 2025b）\footnote{deepseek-ai/DeepSeek-V3.2-Exp-Base}、DeepSeek v4 Pro（DeepSeek-AI, 2026）\footnote{deepseek-ai/DeepSeek-V4-Pro}、Kimi-K2（Kimi Team et al., 2026）\footnote{moonshotai/Kimi-K2-Base} 和 Gemma4-31B（Gemma Team, 2026）\footnote{google/gemma-4-31B}。此外，我们还纳入了我们上一代的一个同样以 30T 数据训练的 23B 模型，以便量化我们的进展。

结果汇总于图 10。MAI-Base-1 在活跃参数数量相近的模型中展现出明显优势，在所有四项留出任务上都取得了更低的每字节比特数。DeepSeek-v4 取得了最佳结果，其活跃参数为 MAI-Base-1 的 1.4$\times$，总参数为 1.6$\times$。我们的


\pandocbounded{\includegraphics[keepaspectratio,alt={image}]{images/bece020ada24a8f0a70718f99308f62049a7a15ecfc154d06f13d563e7f2b599.jpg}}

\pandocbounded{\includegraphics[keepaspectratio,alt={image}]{images/f8fab11bc9951ab19148b2283c4bde7385608977aaecd0eb9a67176a717cef2b.jpg}}

\pandocbounded{\includegraphics[keepaspectratio,alt={image}]{images/15b5f923ac7fb05e425cac69ae239d46febe870db4d2996a9c9d9f2e8e6eb31c.jpg}}

图 10. 基础预训练模型在四项留出评估任务上的每字节比特数（BPB）对比（越低越好）。MAI 模型以黄色标示。我们还额外提供了一个数据点：23B（L66，来自我们 ladder 2.2.1 的旧版本），这是我们此前用与 MAI-Base-1 相同数量的 token（30T）训练的版本。所使用的内部留出评估任务见第 2.7 节。

\pandocbounded{\includegraphics[keepaspectratio,alt={image}]{images/17ae8a898317d9c3305cb6f76f0c670ee6d9fbef6f00b209ae35b6f7f0981e8e.jpg}}

我们的 23B 模型仅使用其 62\% 的活跃参数和相近的总参数量，就超越了 DeepSeek v3.2 基础模型。

\subsection{2.8 YOLO：大规模分布式训练}\label{yolo-distributed-training-at-scale}

You Only Launch Once（YOLO）是 MAI 自研的大规模训练框架，贯穿于 MAI-Thinking-1 的整个开发周期。YOLO 构建于 PyTorch（Paszke et al., 2019）之上，位于第 L.3 节所述的调度与运行时栈内部。

YOLO 旨在支持大规模模型训练的所有阶段，包括预训练、中期训练、监督微调，以及强化学习的训练部分。它实现了核心训练循环，包括我们的模型定义、分片、优化器、数据加载和检查点保存。我们从零构建 YOLO，以对整个训练栈进行端到端控制：从内核到并行再到调度。YOLO 借鉴了 Megatron-Core（Yan et al., 2026）、DeepSpeed（Rasley et al., 2020）和 TorchTitan（Liang et al., 2024）等大规模训练框架的经验，并与我们其余的计算基础设施和模型训练协同设计、深度集成。这使得我们能够在计算环境和模型架构演进过程中，同等优先地考虑训练速度、有效吞吐（goodput）、确定性、灵活性和开发者敏捷性。以下章节描述 YOLO 的主要组件，以及它们如何按照我们的设计原则实现优化。



\subsection{2.8.1 系统概览}\label{system-overview}

YOLO 由一组构建模块组成，用于实现 Transformer 模型，以及训练循环和数据加载器、检查点、日志与可观测性等辅助组件。

内核（Kernels）。为了在大规模下高效支持数值配方（第 2.6.3 节），YOLO 包含一套用 Triton、CUDA、CuteDSL 和 CUTLASS 编写的自定义内核。其中包括针对前向和数据梯度计算优化、并支持 FP8 延迟缩放的 FP8 GEMM 内核；面向 MoE 的分组 GEMM，可在每个 GPU 承载多个专家时提升效率——我们的预训练运行中通常为 8 个（512 个专家，专家并行度 64）——并支持所有混合精度执行模式，包括 FP8 延迟缩放；以及融合量化内核，它对反向传播所需的原始和转置张量同时量化，同时应用缩放因子 swizzling 以提高内存效率，并在可能的情况下将紧邻 RMSNorm 层的转换融合到归一化内核中，以减少开销。此外，非分组量化内核支持集群启动控制（Cluster Launch Control，CLC），这是 Blackwell 特有的用于动态内核负载均衡的功能。

张量分片（Tensor sharding）。根据模型大小和单次训练运行的集群基础设施，YOLO 使用不同形式的并行组合。为了对模型中各张量的分片方式进行细粒度控制，YOLO 实现了一套自定义分片标注，类似于 JAX（Bradbury et al., 2018）或 PyTorch 的 DTensor（Paszke et al., 2019）。分片标注描述了张量如何在 n 维笛卡尔网格上划分或复制，并可用于按分片类型对张量分组。每个张量所划分到的底层网格位于 PyTorch 通信器之上，该通信器为相应张量暴露集合通信和点对点通信原语。我们刻意让标注保持纯描述性，即它们不会自动在计算图中插入通信原语，以避免产生意外的同步点。

YOLO 的分布式张量使我们能够独立定义模型中每个参数的分片方式。这允许我们在模型内部灵活地进行张量并行。例如，我们对嵌入、损失和注意力权重使用张量并行，但对 MLP 和 MoE 权重不使用。另一个例子是，我们可以在属于 MoE 层或非 MoE 层的张量之间使用不同程度的数据并行（称为并行折叠（parallel folding）（Yan et al., 2026））。

并行（Parallelism）。基于分布式张量，YOLO 实现了所有常用的并行形式：数据并行、张量与上下文/序列并行、专家并行和流水线并行（Huang et al., 2019; Qi et al., 2023）。

• 数据并行。YOLO 使用从零构建的自定义版本 ZeRO 阶段 1-3（Rajbhandari et al., 2020），以适配我们的优化器和性能优化。该实现按参数的唯一分片、数据类型和可配置的 ZeRO 单元对参数分组，并将所有参数和梯度合并到连续的缓冲区（buffers）。这最大限度地减少了梯度归约和 ZeRO 1-3 所需的集合通信总数。此外，我们的 ZeRO 实现将参数拆分为多维块，使其与分布式检查点的切片保持兼容，并避免在执行集合通信前后重新排序。无论处于哪个 ZeRO 阶段，YOLO 的 ZeRO 实现始终以分片形式存储参数。这简化了分布式 AdamW 的实现，因为优化器状态也始终自动分片。ZeRO 各阶段的主要区别在于收集参数或梯度的缓冲区何时以及多久被清除。例如，对于 ZeRO-1，参数在第一次前向传播期间全量收集（all-gather），直到最后一个微批次后才清除；而对于 ZeRO-3（FSDP），参数在使用后立即清除。



• YOLO 中的张量并行通过列并行和行并行 GEMM 实现，其中激活和权重张量都在相同数量的 rank 上进行分片（通常称为序列并行）（Shoeybi et al., 2019; Korthikanti et al., 2023）。超出权重分片程度的任何额外激活张量分片，在注意力层中视为上下文并行，在其他情况下视为数据并行。

• 上下文并行。YOLO 使用 Ulysses 风格的上下文并行（Jacobs et al., 2023），其中激活张量通过并行矩阵转置（all-to-all）重新划分，将序列维度的分片转移到注意力头（隐藏）维度。对于 MAI-Base-1，上下文并行仅用于长上下文中期训练，而张量/序列并行用于所有训练阶段。

• 专家并行。最初，YOLO 的专家并行实现类似于 GShard（Lepikhin et al., 2020），使用静态张量形状和预定义容量因子的 all-to-all 通信。在后续运行中，我们切换到了无丢弃（dropless）MoE 实现，通过可变大小的 all-to-all 通信将每个 token 路由到其选定的专家。所有预训练运行都在 NVLink 域内使用专家并行（张量并行度为 1），并在跨机架间使用数据并行。

激活检查点与卸载（Activation checkpointing and offloading）。当使用 ZeRO 数据并行训练 MAI-Base-1 时，激活占用了大部分可用的高带宽设备内存（HBM）。为了减少内存占用并支持在单个 NVLink 域内进行专家并行，YOLO 同时支持激活检查点以及将激活卸载到主机内存。与 Megatron-Core 等框架类似，激活检查点可以按不同粒度配置，从重算一个或多个完整 Transformer 层，到层内中间激活的细粒度重算（例如重算激活函数的输出）。在 MAI-Base-1 的所有训练阶段中，我们结合使用 Transformer 层检查点和细粒度重算。由于存储检查点激活所需的内存量随 Transformer 层数线性增长，YOLO 还实现了激活卸载。激活卸载与检查点互补，且开销极小，因为设备到主机的传输在专用 CUDA 流上使用锁页主机内存异步执行。

专家混合（Mixture of experts）。我们的 MoE 层在研究灵活性与训练效率之间取得平衡。它同时支持无丢弃和容量受限两种模式、多种负载均衡策略（包括全局批次负载均衡损失（Qiu et al., 2025; Yang et al., 2025）和无损失负载均衡（Wang et al., 2024a））、可配置的 top-k 评分、细粒度遥测，以及用于 RL 采样的路由器回放。对于专家并行，我们将本地专家划分为可配置大小的组，并在各组之间流水线化调度 $\to$ 计算 $\to$ 收集各阶段，使得除第一次调度和最后一次收集之外的所有阶段都与专家计算重叠；剩余暴露的通信与全局负载均衡计算以及（可选的）共享专家相重叠。为防止无丢弃模式中因不均衡导致的内存波动、碎片化和 OOM，我们支持静态内存无丢弃模式，每组运行多轮容量受限的调度 $\to$ 计算 $\to$ 收集，每轮最多处理固定容量的 token。在反向传播中，我们执行每专家每轮的细粒度重算，以避免存储不均衡的激活。在性能方面，我们使用低开销地支持所有数值配方的自定义分组 GEMM 和量化内核，以及自定义 CuTe DSL 对称内存内核，用于通过 NVLink 进行设备发起的、可变大小的、高吞吐 all-to-all 通信。

\subsection{2.8.2 确定性与正确性}\label{determinism-and-correctness}

我们强制训练中的确定性，原因包括科学可复现性、开发者敏捷性和系统健康调试，即使以降低训练效率（如 MFU）为代价。对于固定的硬件拓扑、模型配置和软件版本，我们确保预训练过程达到逐位可复现。简而言之，使用相同硬件资源和配置的两次训练运行将产生逐位一致的模型。确定性需要在训练栈的多个层面进行仔细处理：



• 数据管线。对于给定的数据并行 worker 配置，训练微批次的顺序是固定的。虽然这些 worker 会异步预取批次以隐藏 I/O 延迟，我们确保待处理批次的队列在多次运行和重启之间保持一致。

• 训练检查点。我们存储所有有状态数据，包括模型权重、优化器状态、FP8 缩放历史、数据加载器进度和随机数生成器。

• GPU 内核。浮点累加不可结合，因此可复现的归约需要控制值的合并顺序。我们的训练栈按规定顺序执行确定性累加，而不是依赖无序的 GPU 原子操作（其累加顺序取决于调度器）。例如，RMSNorm 反向传播使用两阶段分块归约累加梯度：第一个内核物化每个分块的部分和，一个最终化内核按固定顺序归约这些部分和。另一个例子是 MoE 路由的 top-k 选择：我们的内核使用稳定排序以避免非确定性的平局打破。

• 网络集合通信。与 GPU 内核内的浮点归约类似，all-reduce 和 reduce-scatter 等集合通信必须是确定性的。我们禁用 NVLink SHARP（Graham et al., 2020），以确保集合操作确定，代价是性能下降。我们还强制每次运行采用一致的 NCCL 拓扑，使机架内归约始终使用相同的累加顺序。

我们还利用确定性来建立对 YOLO 演进过程中数值正确性的信心。我们的测试跟踪 ``golden'' 模型配置，并与所有训练指标（如损失、梯度范数和权重张量校验和）进行精确性对比。这并不意味着我们禁止任何会导致不同训练结果的数值改动；当我们做出影响浮点行为（例如升级系统软件或 GPU 内核）的改动时，出现变化是预期之中的。这为我们提供了软件历史中预期数值变化的明确记录。

\subsection{2.8.3 容错}\label{fault-tolerance}

大规模训练必须容忍频繁的硬件故障，并将对训练效率的影响降至最低。YOLO 的容错方法结合了分布式检查点和向热备（hot standby）的快速故障转移，以最大限度地减少故障期间的训练进度损失和停机时间。

分布式检查点格式（Distributed checkpoint format）。我们基于 PyTorch 的分布式检查点库（DCP）构建，该库对跨 rank 的复制状态进行去重，并支持加载时的张量重分片。我们重写了 DCP 的元数据创建和序列化路径，以降低 CPU 内存开销、减少 Python 垃圾回收压力，并加速序列化/反序列化。总体而言，这些优化将 CPU 内存开销和检查点保存时间都降低了 10$\times$ 以上。在我们的实现中，原库仅保留了公共 API 和将张量分片映射到存储的核心逻辑——即保存计划（save plan）。我们还预先计算 DCP 保存计划并缓存在持久化存储中，将该昂贵计算从关键路径中移除。

异步检查点（Asynchronous checkpointing）。为最大限度地减少保存检查点时的训练停机，我们首先在训练进程中直接将检查点张量从设备复制到主机，然后在训练继续的同时将剩余工作卸载到一个单独的检查点进程。检查点进程通过前端网络将数据写入持久化存储，并处理跨 rank 协调以原子地提交检查点。为支撑我们的规模，我们用构建在共享键值存储之上的自定义协调层替换了标准 CPU 集合通信后端。我们还采用机会式准入控制机制，轮询进行中的检查点是否完成，并尽快启动下一个检查点，同时确保任意时刻最多只有一个检查点在进行中。这种方法与细致的主机内存管理相结合，使我们能够频繁保存检查点，同时将对训练 MFU 的影响降至最低。



存储基础设施（Storage infrastructure）。我们的训练数据中心拥有本地 Azure Blob Storage（Calder et al., 2011），提供低延迟、高带宽的访问。这使我们能够直接依赖检查点实现容错，而无需中间节点本地缓存或跨 rank 状态复制。为避免检查点加载时大规模出现存储热点（许多 rank 需要相同的复制状态），我们仅从存储加载一次复制状态，并通过 NCCL 广播，消除了可能压垮单个存储服务器的单 blob 扇入热点。

快速任务恢复（Rapid job recovery）。我们通过使用 Ray（Moritz et al., 2018）actor 配合热备实现任务内重启，从而获得快速恢复，省去了 pod 重建的开销。在 GB200 系统上，我们在 actor 重启期间显式清除操作系统缓存，以防止因 GPU 持有上次运行的残留内存导致内存不足错误而引发的崩溃循环。最后，我们通过将重算步骤与历史损失指标对比，在重启后验证确定性。

\subsection{2.8.4 与模型架构协同优化性能}\label{cooptimizing-performance-with-model-architecture}

MAI-Base-1 的最终架构是跨越模型五个主要世代（v1--v5）的持续协同设计过程的成果。每一世代追求两个互补的目标。首先，我们通过改变模型设计、路由策略和训练配方，力求提升架构的效率增益与建模能力。其次，我们通过识别性能瓶颈并引入新的基础设施、内核、通信和内存管理优化，持续提升训练效率。由于架构改动改变了模型的计算特性，它们常常引入新的系统瓶颈，需要在训练栈中做出相应的改进。

我们描述模型的架构与基础设施演进，以突出我们构建大规模模型的方法。每次演进中跟踪的两个指标是：

• 模型 FLOP 利用率（Model FLOP Utilization，MFU）（Chowdhery et al., 2023），定义为
\(\mathrm { F L O P } / ( t _ { \mathrm { s t e p } } { \cdot } \mathrm { F L O P } _ { \mathrm { s p e c } } )\)
用于量化实际训练效率。这里 FLOP 表示所有前导阶计算内核的浮点运算总数，包括 GEMM 和注意力，但不包括 RMSNorm 和激活函数等内存受限操作。
\(t _ { \mathrm { s t e p } }\) 是端到端训练步时间，包括数据加载、前向和反向传播、通信以及优化器更新。尽管训练混合使用 FP8、BF16 和 FP32 运算，我们使用 NVIDIA 报告的 FP16/BF16 张量核心吞吐量来归一化 MFU（FLOPsspec = 2.5 $\times$ 1015 FLOPS 每 GB200 GPU，用于稠密运算）。激活检查点引入的运算不计入 FLOP 计数。

• 效率增益（Efficiency Gain）：定义见第 2.2.2 节，用于刻画在达到给定质量时模型效率的理论提升。

表 7 汇总了各模型世代间的关键架构改动。图 11 同时报告了新架构配合上一世代优化栈运行时的初始 MFU，以及应用针对各版本的性能改进后达到的最终 MFU。在 MAI-Base-1 的开发过程中，引入了超过二十项基础设施和内核优化，使每一世代在模型规模和架构复杂性大幅增加的情况下仍能维持 20\% 以上的 MFU。世代 v0 和 v1 是在 H100 系统上训练的早期原型，因此未包含在图 11 中。

从版本 2 开始，所有主要训练运行都在 GB200 NVL72 集群上进行。MAI-Base-1 的版本 2（23B 活跃参数）是首个在 GB200 NVL72 集群上训练的模型版本，共使用 4,096 块 GPU。为提高系统鲁棒性，我们每机架使用 64 块 GPU（NVL64），预留冗余容量以容忍节点故障和不健康设备。版本 2 采用了比早期版本更深更窄的设计，同时保持 23B 活跃参数预算，因此隐藏维度相对较小。为最大化 GEMM 效率，我们选择了专家并行（EP）度为 64、张量并行（TP）度为 1，确保所有专家 all-to-all（A2A）通信都保持在 NVL64 域内。该配置实现了 18\% 的初始 MFU。我们使用跨机架 InfiniBand（IB）进行数据并行（DP）通信（如参数 all-gather 和梯度 reduce-scatter）。我们在后续版本中保持了相同的设置。



表 7. MAI-Base-1 开发过程中的架构演进。这里报告的是用于大规模基础设施测试的目标架构，而非 ladder 模型。

{\def\LTcaptype{none} % do not increment counter
\begin{longtable}[]{@{}|l|l|l|l|l|@{}}
\toprule\noalign{}
\endhead
\bottomrule\noalign{}
\endlastfoot
\hline
模型版本 & v2 & v3 & v4 & v5 (MAI-Base-1) \\
\hline
激活/总参数量 & 23B/600B & 23B/600B & 23B/611B & 35B/1T \\
\hline
层数 & 54 & 54 & 66 & 78 \\
\hline
隐藏维度 & 4352 & 4352 & 5632 & 6656 \\
\hline
FFN & 34816 & 34816 & 11264 & 13312 \\
\hline
Down-proj 维度 & & & 2560 & 3072 \\
\hline
专家 FFN & 17408 & 17408 & 9216 & 10240 \\
\hline
专家容量因子 & 2 & $\infty$ & $\infty$ & $\infty$ \\
\hline
Top-k/专家数 & 4/192 & 4/192 & 8/512 & 8/512 \\
\hline
\end{longtable}
}

性能剖析揭示了三个主要瓶颈：DP 通信、确定性注意力内核和专家路由开销。启用 GPU Direct RDMA 支持改善了通信与计算的重叠，并将端到端训练时间缩短了约 1.1$\times$。由于 FlashAttention-2（FA2）的确定性模式在 GB200 上效率低下，且 FA4 不可用，我们开发了一个自定义块稀疏注意力后端，将步时间额外缩短了 1.06$\times$。鉴于单微批次训练配置，我们采用 ZeRO-2 优化器分片，与 ZeRO-1 相比，在不引入额外通信的情况下降低了梯度内存消耗。这使我们能够禁用额外的激活重算，并进一步提升 1.03$\times$ 的性能。最后，我们将一个内存受限的 PyTorch 专家编码内核（仅利用 10\% 的可用 HBM 带宽）替换为自定义 Triton 实现，达到了约 80\% 的 HBM 利用率，并将端到端性能额外提升了 1.03$\times$。综合起来，这些优化将 MFU 从 18\% 提升到 22\%。

版本 3 保持了相同的整体模型架构，但从容量受限路由切换为无丢弃（dropless）MoE 路由。从效率增益的角度看，这使我们能够采用更好的 MoE 负载均衡策略。这对 MFU 也有积极影响，因为消除了专家容量填充，从而同时减少了通信量和专家 GEMM 计算量。然而，无丢弃路由引入了新的同步开销，因为专家 token 数和张量形状必须在执行期间动态通信。为缓解这些成本，我们将 token 数通信与其他运行时操作重叠，并将同步点移至专用流上。这些改动使我们在保持与版本 2 相当 MFU 的同时，实现了无丢弃路由的效率收益。

版本 4 引入了三项重要的架构改动：专家数量从 192 增加到 512，路由从 top-4 扩展到 top-8 专家选择，并使用了 LatentMoEs。增加专家数量有助于提升模型的学习容量，而 latentMoE 有助于在专家数量增加时管理 A2A 通信成本。同时，训练任务从 4,096 块 GPU 扩展到 8,192 块 GPU。LatentMoEs 减小了专家 GEMM 的规模，使得 CPU 启动开销和内核效率变得更加重要。结果是，初始 MFU 从 22\% 下降到约 16\%。使用 FA4 确定性内核（相对于块稀疏注意力后端提升 1.14$\times$），再加上若干旨在降低 CPU 开销、提升运行时批处理效率的改进，将 MFU 从 16\% 提升到 20\%。

版本 5 将活跃参数量从 23B 增加到 35B，总参数量从 600B 增加到 1T。更大的模型显著增加了参数和激活内存需求。因此，初期部署依赖 ZeRO-3 优化器分片，但性能剖析显示，额外的参数 all-gather 操作导致反向传播受通信限制。为解决这一问题，我们通过激活卸载降低 GPU 内存压力，使训练系统得以回归 ZeRO-2 分片。消除 ZeRO-3 参数 all-gather 操作恢复了通信与计算之间的重叠，并提升了整体训练效率，使最终版本 5 配置能够维持约 20\% 的 MFU。



图 11. 在 GB200 上训练的各类预训练配置中 MFU 和 EG 的演进。从 v2 起，每次模型改动都提升了 EG，但在部署优化前用上一代配置运行时，MFU 起初会下降。总体上，我们为每次预训练运行添加了 20+ 项优化，以超越 20\% 的 MFU。注意，这里我们只列出具有显著基础设施影响的模型改动——它们并非模型各版本之间促进 EG 提升的全部改动。

\pandocbounded{\includegraphics[keepaspectratio,alt={image}]{images/3555fbd21605b657f2f72c8780a3f004b523efd7e02488948b7aed4e3eef13d8.jpg}}
